\section{Contribution and Expected Outcomes}
\label{sec:contribution}

\subsection{Contribution of the Study}

This study contributes to the field of emotion detection by providing a structured comparative evaluation of traditional machine learning and deep learning approaches under a unified experimental framework. Unlike studies that evaluate models under varying preprocessing conditions or heterogeneous datasets, this research implements a consistent NLP pipeline, standardised dataset partitioning strategy, and uniform evaluation metrics to ensure fair and interpretable comparison.

The study further contributes by analysing the influence of feature representation methods --- including TF-IDF and embedding-based approaches --- within the same dataset context. Through class-wise performance evaluation and confusion matrix analysis, the research provides deeper insight into model behaviour across discrete emotion categories. Concretely, this study will produce a reproducible experimental benchmark --- including preprocessing scripts, model training code, and a performance comparison table --- made available for future researchers to directly extend or compare against on the same 416,809-sample Twitter corpus.

\subsection{Expected Outcomes}

Deep learning models, particularly BiLSTM, are expected to outperform traditional ML algorithms due to their ability to capture contextual and sequential dependencies. CNN models are anticipated to demonstrate competitive performance through effective local feature extraction, while SVM and XGBoost are expected to provide reliable baseline results with lower computational complexity. Certain emotion pairs --- particularly fear/surprise and sadness/love --- are expected to yield higher misclassification rates due to semantic similarity, providing meaningful insight into the limits of each architectural approach.
